%%%%%%%%%%%%%%%%%%%%%%%%%%%%%%%%%%%%%%%%%
% Tufte Essay
% LaTeX Template
% Version 2.0 (19/1/19)
%
% This template originates from:
% http://www.LaTeXTemplates.com
%
% Authors:
% The Tufte-LaTeX Developers (https://www.ctan.org/pkg/tufte-latex)
% Vel (vel@LaTeXTemplates.com)
%
% License:
% Apache License, version 2.0
%
%%%%%%%%%%%%%%%%%%%%%%%%%%%%%%%%%%%%%%%%%

%----------------------------------------------------------------------------------------
%	PACKAGES AND OTHER DOCUMENT CONFIGURATIONS
%----------------------------------------------------------------------------------------

\documentclass[a4paper]{tufte-handout} % Use A4 paper by default, remove 'a4paper' for US letter

\usepackage{graphicx} % Required for including images
\setkeys{Gin}{width=\linewidth, totalheight=\textheight, keepaspectratio} % Default images settings
\graphicspath{{Figures/}{./}} % Specifies where to look for included images (trailing slash required)

\usepackage{booktabs} % Required for better horizontal rules in tables

\usepackage{subfiles}

%----------------------------------------------------------------------------------------
%	TITLE SECTION
%----------------------------------------------------------------------------------------

\title{Algebraic and Transcendental functions}

\author{Professor: Francisco Vazquez-Tavares}

\date{Last compilation: \today} % Date, use \date{} for no date


%----------------------------------------------------------------------------------------
%	COMMANDS SECTION
%----------------------------------------------------------------------------------------


%----------------------------------------------------------------------------------------

\begin{document}

\maketitle % Print the title section

%----------------------------------------------------------------------------------------
%	ABSTRACT/SUMMARY
%----------------------------------------------------------------------------------------

\begin{abstract}
    \textbf{\large Name: }
\end{abstract}

%----------------------------------------------------------------------------------------
%	ESSAY BODY
%----------------------------------------------------------------------------------------


\marginpar{
    {\bfseries\Large Algebra Toolkit}

    {\bfseries\large Algebra properties}
    \begin{gather*}
        a+b = b+a, \\
        \qty(a+b)+c = a+\qty(b+c), \\
        a\qty(b+c) = ab+ac, \\
        \frac{a}{b} + \frac{c}{d} = \frac{ad+bc}{bd}, \quad
        \frac{a}{b} \cdot \frac{c}{d} = \frac{ac}{bd}, \\
        \frac{-a}{b} = -\frac{a}{b} = \frac{a}{-b}, \\
        \frac{\frac{a}{b}}{\frac{c}{d}} = \frac{ad}{bc}
        \\
        \sqrt{ab} = \sqrt{a}\sqrt{b},\quad
        \sqrt{\frac{a}{b}} = \frac{\sqrt{a}}{\sqrt{b}},\quad
        \sqrt[n]{a} = a^{1/n}
        \\
        x=\frac{-b\pm\sqrt{b^2-4ac}}{2a},\\
        a^2-b^2 = \qty(a-b)\qty(a+b)
    \end{gather*}
    
    {\bfseries\large Laws of exponents}
    \begin{align*}
        a^r\cdot a^s = a^{r+s}, \quad & \qty(a^r)^s = a^{rs} \\
        \frac{a^r}{a^s} = a^{r-s}, \quad & \qty(ab)^r = a^r b^r \\
        \qty(\frac{a}{b})^r = \frac{a^r}{b^r},\quad & b\neq 0
    \end{align*}

    {\bfseries\large Laws of Logarithms}
    \begin{align*}
        \log_a\qty(AB) &= \log_a\qty(A) + \log_a\qty(B) \\
        \log_a\qty(\frac{A}{B}) &= \log_a\qty(A) - \log_a\qty(B) \\
        \log_a\qty(A^C) &= C\log_a\qty(A)
    \end{align*}

    {\bfseries\large Inverse functions property}
    \begin{gather*}
        \qty(f\circ f^{-1})(x) = x = \qty(f^{-1}\circ f)(x) \\
        \log_a(a^x) = x = a^{\log_a(x)}
    \end{gather*}
}

\section{Algebra skills}

Combine the functions $f(x) = x^2+1$ and $g(x)=x-4$, according to the indicated operation,
\begin{itemize}
    \item $(f+g)(x)$
    \item $\left(\frac{f}{g}\right)\left(0\right)$
\end{itemize}

Solve the following equations,
\begin{itemize}
    \item $\log_2\qty(2x+2) - \log_2\qty(x+1) = 3$
    \item $2\log_3\qty(x) - \log_3\qty(x-4) = 2+\log_3\qty(2)$
\end{itemize}

\section{Function analysis}

For each of the following pair of function, find the indicated composition, its domain and range.

\begin{itemize}
    \item $f(x) = x^2 + 7$ and $g(x)=\sqrt{x+5}$
    \item $f(x) = \frac{7x}{x+2}$ and $g(x) = \frac{5}{x}$
\end{itemize}

Find all the asymptotes, $x$ and $y$ interceptions of the following function,
\begin{gather*}
    f(x) = \frac{2x-4}{2x^2+x-10}
\end{gather*}


\section{Modelling with math}

\paragraph{Publishing company.} A small publishing company is releasing a new book.
The production costs will include a one-time fixed cost for editing and an additional corst for each book printed.
The total production cost $C$ (in dollars) is given by the function $C=650+15.9N$, where $N$ is the number of books.
The total revenue earned (in dollars) from selling the books is given by the function $R=32.80N$.
Let $P$ be the profit made (in dollars). 
\begin{itemize}
    \item Write an equation relating $P$ to $N$. 
    \item Simplify your answer as much as possible.
\end{itemize}

\paragraph{Population growth.} The population of a certain city is given by $P(t)=234e^{0.071 t}$, where $P$ is measured in thousands of people and $t$ is measured in years from the year 2000.
\begin{itemize}
    \item Find the population for the year 2014.
    \item Find the population for the year 2027.
        find when the population will be 3 million.
\end{itemize}

\newpage

\marginpar{
    {\bfseries\Large Conic Sections}

    {\large Standard form of a circle}
    \begin{gather*}
        (x-h)^2 + (y-k)^2 = r^2
    \end{gather*}

    {\large Standard form of a parabola}
    \begin{gather*}
        (x-h) = 4a(y-k)^2,~\mathrm{or}~(y-k) = 4a(x-h)^2
    \end{gather*}

    {\large Standard form of an ellipse}
    \begin{gather*}
        \frac{(x-h)^2}{a^2} + \frac{(y-k)^2}{b^2} =1
    \end{gather*}

    {\large Standard form of a hyperbola}
    \begin{gather*}
        \frac{(x-h)^2}{a^2} - \frac{(y-k)^2}{b^2} =1
    \end{gather*}

}

\section{Conic sections}

Identifty the type of conic section and find the following,

\vspace{2em}

\begin{minipage}[c]{0.45\textwidth}
\begin{itemize}
    \item If is a cirlce,
        \begin{enumerate}
            \item Center
            \item Radius
        \end{enumerate}
    \item If is parabola
        \begin{enumerate}
            \item Vertex
            \item Focus
            \item Directrix
        \end{enumerate}
\end{itemize}
\end{minipage}
\begin{minipage}[c]{0.45\textwidth}
\begin{itemize}
    \item If is an ellipse
        \begin{enumerate}
            \item Center 
            \item Focus (2)
            \item Semi-axis (2)
        \end{enumerate}
    \item If is a hyperbola
        \begin{enumerate}
            \item Center
            \item Focus (2)
            \item Asymptotes (2)
        \end{enumerate}
\end{itemize}


\end{minipage}



\begin{enumerate}
    \item $x^2 -4x +4 + y^2 + 6y +9 = 9$
    \item $x^2 + (y-2)^2 = 36$
\end{enumerate}



\end{document}
